% Mayo 2020 (modificada por Juande Santander-Vela, profesor BBDD y Observatorio Virtual)
% Ejemplo básico LaTeX (basado en la plantilla original de Rafael  Padilla,
% alumno de la primera edición % del máster de Astronomía y Astrofísica de la VIU)
% Lo primero: Las líneas precedidas por "%" son de comentarios y el compilador LaTeX las ignorar.
% Se ponen para dar más claridad a lo que se está haciendo. A continuación, el código en sí.
% Los comandos de LaTeX comienzan por una barra invertida ( \ ).


% Preámbulo de LaTeX. %%%%%%%%%%%%%%%%%%%%%%%%%%%%%%%%%%%%%%
% Definición del tipo de documento.
\documentclass[a4paper]{article}


% Paquetes que usaremos:
% Use UTF8 as the deafult character encoding. Works better for most modern systems.
\usepackage [utf8]{inputenc}

\usepackage[spanish,activeacute]{babel}


% Graficos
\usepackage{graphicx}

% Este nos activa el paquete "babel" con sus acentos idioma y guionado en español.
%\usepackage[spanish,activeacute]{babel}

%Este otro activa el uso de cierta nomenclatura y modo de representar para ecuaciones matemáticas.
\usepackage{amsmath}

% BibTeX
\usepackage[sectionbib]{natbib}  % Cross-reference package (Natural BiB)
\bibpunct{(}{)}{;}{a}{}{,}

% Títulos del documento.
\title{Introducción al \LaTeX{} para alumnos de la VIU}
\author{Autor original: Rafael Padilla}
%\date{2021-01-01} % Usa \date para cambiar la fecha a mostrar; comenta para mostrar la fecha de hoy




% Fin del Preámbulo. %%%%%%%%%%%%%%%%%%%%%%%%%%%%%%%%%%%%%%%%


% Contenido. %%%%%%%%%%%%%%%%%%%%%%%%%%%%%%%%%%%%%%%%%%%%
\begin{document}

% Esto nos pone los títulos de acuerdo al tipo de documento definido.
\maketitle

% Descomenta los siguientes comandos si quieres tener tabla de contenido, lista de figuras, y/o lista de tablas
%\tableofcontents
%\listoffigures
%\listoftables

\begin{figure}[!t]
\begin{center}
\includegraphics[height=2cm]{figures/logo_viu}
\end{center}
\end{figure}


\section{Introducción}
Nuevamente me atrevo a escribir un pequeño texto, esta vez sobre~\LaTeX{}, para tratar de siquiera aportar un granillo de arena que sirva, no para que se nos meta en el ojo, sino para sacarlo de el. Para que nos ayude —a mi el primero— a entender bien este asunto del \LaTeX{}, el cual de verdad es una virguería\footnote{Para aquellos no acostumbrados a léxico similar por no ser utilizado en su ubicación, diremos que virguería es sinónimo de cosa ciertamente habilidosa, poderosa, a veces sobrecogedora y siempre, siempre, bastante versátil.}.

\subsection{¿WYSIWYG?}
Lo primero, sería preguntarse: ¿Cómo un lenguaje basado en comandos de línea ha perdurado tantísimo tiempo, sin ser reemplazado por cosas del estilo WYSIWYG y procesadores de texto llenos de ventanas? Hay varias razones. La primera, por su potencia innegable e insuperable a la hora de publicar textos científicos. Todo el mundo sabe que el formato textual del Astrophysical Journal no es el mismo que el del Muy Interesante, por poner un ejemplillo. La primera es una revista austera, potente, seria. Muy científica y con un formato conocido a la hora de publicar datos. La ciencia no ha parecido, por lo general, necesitar demasiados colorines ni juegos tipográficos.

El bosque tropical desaparece a enorme velocidad:
\begin{enumerate}
\item Cada 2 segundos se tumba un área de bosque gigantesca.
\item Cada día se destruye un área del tamaño de la Ciudad de Guatemala.
\end{enumerate}

El bosque tropical desaparece a enorme velocidad:
\begin{itemize}
\item Cada 2 segundos se tumba un área de bosque gigantesca.
\item Cada día se destruye un área del tamaño de la Ciudad de Guatemala.
\end{itemize}

Por otra parte, podría decirse que, dado que los artículos son escritos directamente por los propios científicos que publican utilizando las plantillas de \LaTeX{} de ApJ\footnote{Que como todos saben, significa... ¿Qué era, eh?}, pues la verdad, se obtiene un estilo uniforme, directo, científico y austero. \LaTeX{} permite mucha virguería gráfica a la hora de representar datos, figuras, gráficas y demás; existen módulos de \LaTeX{} para las partidas de ajedrez, para la música, para el juego del Go... Muchas, muchísimas publicaciones se realizan en este soberbio lenguaje para la creación profesional de documentos.

\LaTeX{} posee una ventaja adicional una vez dominado: Cuando todos los elementos están dispuestos y listos, cuando se han dispuestos los comandos de la forma adecuada\footnote{Y para eso están las plantillas, precisamente: para ahorrar mucha pejiguera al redactor!}, el autor del texto tan sólo ha de preocuparse de una sola cosa: del contenido. Y de nada más. A la hora de enviar éste \emph{a donde sea}, lo que se envía es simplemente texto en ASCII, liviano, ligero, enviable por email. Un documento de 100 páginas puede ocupar unos 250 kB en \LaTeX{}, mientras que si el mismo documento se redacta o se  hace en formato .pdf o en otro de los formatos habituales para texto, éste llegará a ocupar fácilmente varios MB.



\begin{figure}[ht]
\resizebox{\hsize}{!}{\includegraphics{figures/modos}}
\caption{Diferentes modos presentes en una estrella pulsante.}
\label{fig:modos}
\end{figure}


Bueno, ya hemos visto dos ventajas de todo esto: La primera, que cuando está bien configurado el documento, el autor sólo tiene que preocuparse de teclear lo que sea y nada de dar formato al texto, de contar las páginas, de ver que esta figura no cabe y no se sabe donde cortar el texto o como representar estos datos rarísimos en 3D de forma profesional. No. El autor teclea y se ocupa tan sólo \emph{de lo suyo}. Esto es vital, radical y ha convertido a los universitarios en grandes publicadores profesionales. Hay vida más allá de Word.\\

\begin{figure}[ht]
\begin{center}
\includegraphics[width=5cm]{figures/miktex.eps}
\caption{Un sistema TeX{} de 32 bits.}
\end{center}
\label{fig:miktex}
\end{figure}

Pero he aquí que de repente, sin saber ni cómo y sin comerlo ni beberlo, a nosotros nos ha caído encima tantísimo bien con esto del \LaTeX{}. A nosotros, sí. Y por ello escribo este documento. Como en muchas otras cosas, lo que es verdaderamente bueno, muchas veces tiene un precio que pagar por adelantado. En nuestro caso, ese precio puede llamarse perfectamente\footnote{Y con toda la razón del mundo además; no discutiré eso para nada.} \emph{Curva de Aprendizaje}.\\

áEs cierto! Sería muy poco razonable el negar que el que quiere aprender \LaTeX{} tiene que aprender algo de ello. Y en nuestro caso, tenemos la suerte de que esa curva no es demasiado pronunciada; antes bien, lo que tenemos que aprender es como manejarnos por las plantillas y algún que otro de los rudimentos de este software. Para demostrar lo que digo, veremos como está configurado un documento simple, como este.\\

Antes de eso, sin embargo (un poco de paciencia, que enseguida vamos al grano), desearía recomendar una página \emph{fundamental}para todos los TFMistas de la VIU. Esta página es http://www.cervantex.es.\\

\section{Fórmulas matemáticas sencillas}

En lo que sigue $f(x)=\frac{1 + x}{x}$, mientras que \( h(x)= f(a+x)\).


En lo que sigue $$f(x)=\frac{1 + x}{x},$$ mientras que \[ h(x)=f(a+x) \]

Exponentes y radicandos:

$$xy +x^{2z+1}- x^{w-1}_{2}+x^{(z+1)^w} + (x_1 + x_2 )^2$$

$$\sqrt{x-y+z} -\sqrt[4]{(x^2+y^2)^3}$$

Símbolos matemáticos:

$$ a\in A, b\not\in B, A\subset B\not\subset C $$

$$
\lim_{n \to \infty} \sum_{k=1}^n \frac{1}{k^2} = \frac{\pi^2}{6}
$$



\section{Entornos de tablas}

\begin{table}[ht]
\centering
\begin{sffamily}
\begin{tabular}{|llllll|}
\hline
 bahasa	& breton		& catalan 		& croatian 	& czech 		& danish   \\
 dutch	& english		& esperant 	&estonian 	&finnish 		& frenchb\\
 galician 	& germanb	& ngermanb 	& greek		& athnum 		& grmath\\
 irish 	& italian 		&lsorbian 		& magyar 		& norsk 		& polish\\
 portuges	& romanian 	& russianb 	&scottish 		& slovak 		& slovene\\
 spanish 	& swedish 	& turkish 		&ukraineb		& usorbian 	&welsh \\
\hline
\end{tabular}
\end{sffamily}
\label{tab:OpcBabel}
\caption{Opciones del paquete babel. Las opciones athnum y  grmath son extensiones de greek.}
\end{table}


\section{Ecuaciones más complicadas}

Un ejemplo de ``array'':

\begin{equation}
\label{ec:solucion}
\begin{array}{lcl}
\xi_{r}(r,\theta,\varphi,t) 		& = & a(r)Y_{l}^{m}(\theta,\varphi)\exp(i\ 2\pi\nu t)\\
\xi_{\theta}(r,\theta,\varphi,t) 	& = & b(r)\frac{\partial Y_{l}^{m}(\theta,\varphi)}{\partial\theta}\exp(i\ 2\pi\nu t)\\
\xi_{\varphi}(r,\theta,\varphi,t) & = & \frac{b(r)}{\sin \theta}\frac{\partial Y_{l}^{m}(\theta,\varphi)}{\partial\varphi}\exp(i\ 2\pi\nu t)
\end{array}
\end{equation}


where $\xi_{r}$, $\xi_{\theta}$ and $\xi_{\varphi}$ are the displacements, $a(r)$ and $b(r)$ are amplitudes, $\nu$ is the oscillation frequency and $Y_{l}^{m}(\theta,\varphi)$ are the spherical harmonics given by


\begin{equation}
Y_{l}^{m}(\theta,\varphi)=\sqrt{\frac{2\ l + 1}{4\pi}\frac{(l-m)!}{(l+m)!}}\ P_{l}^{m}(\cos\theta)\exp(i\ m\ \varphi)
\end{equation}


that represents the dependence of the mode on the angular variables $\theta$ and $\varphi$ for a star with a spherically symmetric equilibrium configuration. And $P_{l}^{m}(\cos\theta)$ are Legendre polynomials given by


\begin{equation}
P_{l}^{m}(\cos\theta)=\frac{(-1)^{m}}{2^{l}\ l!}\ (1- \cos^{2}\theta)^{\frac{m}{2}}\ \frac{d^{l+m}}{d\ \cos^{l+m}\ \theta}\ (\cos^{2}\theta - 1)^{l}
\end{equation}


where $\theta$ is measured from the pulsational pole, the axis of symmetry. In most pulsating stars that axis coincides with the rotation axis. The main exceptions are the rapidly oscillating Ap stars where the axis of pulsational symmetry is the magnetic axis which is inclined to the rotational axis.

And here we cite the work of~\citet{brown61}, to show how references work.

\section{Conclusión}
En ella está todo lo que podemos necesitar. Por ejemplo, la famosa Introducción no-tan-corta al \LaTeX{}. Está en la sección de Manuales. Mirad. Os recomiendo que echéis una ojeada a la FAQ española y después a la citada guía. No voy a recomendar más por ahora. Seré más extenso de forma específica en mi mensaje siguiente, en el que hablaré de las partes de un documento \LaTeX{} y de algunos de los comandos. Mientras tanto,  disfrutad bien de las guías de Cervantex. En nada estaremos escribiendo todos como papagayos. Y en el mundo nuestro, es una herramienta imprescindible.

Con mis mejores deseos. Saludos,\\
Rafa Padilla.

%% CITAS BIBLIOGRÁFICAS (BibTeX y thebibliography:
%\vspace{1cm}
%\noindent P.D.: Finalmente me gustaría citar a \cite{tal1} y \cite[aunque nunca lo acabé de leer]{bro}.
%




%% thebibliography
%% espacio para el ítem más largo {CLMPS}
%% con 5 caracteres
%
%\begin{thebibliography}{CLMPS} 
%
%\bibitem[Car97]{tal1}
%\textsc{Carlisle, D.}: ``The showkeys package'', documentación distribuida 
%junto al paquete, 1997. 
%
%\bibitem[CLMPS]{bro}
%\textsc{Cascales, B., Lucas, P., Mira, J.~M., Pallarés, A. y 
%Sánchez-Pedreño, S.}: \emph{\LaTeX: una imprenta en sus manos}, ADI, 2000 
%
%\bibitem[Knu86b]{Metafont}
%\textsc{Knuth, D.~E.}: \emph{The MetaFont book}, Addison-Wesley, 1986. 
%
%\bibitem[Knu86a]{TexBook}
%\textsc{Knuth, D.~E.}: \emph{The {\TeX}book\/}, Addison-Wesley, 1986. 
%
%\bibitem[Lam94]{Leslie}
%\textsc{Lamport, L.}: \emph{\LaTeX --- A document Preparation System. 
%User's guide and manual references}, Addison-Wesley, 2nd edici\'on, 1994. 
%
%\bibitem[Spi90]{Spivak}
%\textsc{Spivak, M.~D.}: \emph{The Joy of \TeX. A Gourmet Guide to 
%Typesetting with the AmS\TeX\ macro package}, American Mathematical 
%Society, 1990. 
%
%\end{thebibliography}




% Bibtex:
\bibliography{mibiblio}
\bibliographystyle{sty/astron} % astron para formato de astronomía (incluir el estilo natbib en el preámbulo)


% El comando que sigue marca el final del documento.
\end{document}

